\documentclass{article}
\usepackage{zref-clever}
\usepackage{../latex-course}
% \coursename{Introduction to \LaTeX{}}
\courseweek{02}{Structure, References, Errors}
\coursestudents{Sandrine Pudukkattukaran Joy \and Winter Felix}

\title{Debugging File}
\author{Sandrine Pudukkattukaran Joy (4743017), Winter Felix (4751930)}
\begin{document}
\maketitle

\section{Short Exposition on Tensor Products}\label{sec:tensors}
Throughout this document, let $U$, $V$ and $W$ be vector spaces over one common field $K$.

One way to motivate tensor products is the wish to study \emph{bilinearity} and specifically bilinear maps $b \colon U \times V \to W$ using the tools already developed for \emph{linearity}.
We look for a single vector space $\mathcal{T}$, the \textbf{tensor product space} of $U$ and $V$, together with one fixed bilinear map
\[
	\otimes \colon U \times V \to \mathcal{T},
\]
through which every bilinear map can be factored.
This can be described by the so-called \emph{universal property}, which we will discuss in the following \zcref{subsec:universal}.

\subsection{Definition via the Universal Property}\label{subsec:universal}
The pair $(\mathcal{T}, \otimes)$ is called a \textbf{tensor product} of $U$ and $V$ if it has the following \emph{universal property}: for every bilinear map $b \colon U \times V \to W$ there exists a \emph{unique} linear map $f \colon \mathcal{T} \to W$ with $b = f \circ \otimes$, that is,
\[
	b(u,v) = f(u \otimes v)
\]
for all $u \in U$ and $v \in V$.
Investigating, we can find that mapping a $b$ to its linear partner $f$ is a vector space isomorphism
\[
	\mathrm{Bil}(U,V;W) \; \cong \; \mathrm{Hom}(\mathcal{T}, W).
\]
Remarkably, this single property already fixes $\mathcal{T}$ uniquely up to isomorphisms in a tensor product sense.
Obviously, this definition does not immediately guarantee existence of a tensor space.
Examples can be constructed using bases or as suitable quotient spaces.

\subsection{Tensors, Elementary Tensors and Bilinearity}\label{subsec:tensors}
The elements of $\mathcal{T}$ are called \textbf{tensors}, and the zero vector of $\mathcal{T}$ is the \emph{zero tensor}.
The images of the universal map, that is the tensors of the form $u \otimes v$ with $u \in U$ and $v \in V$, are the \textbf{elementary tensors} (also called \emph{simple tensors}); one reads $u \otimes v$ as ``$u$ tensor $v$''.
Be careful: Tensors are not necessarily simple tensors, but they can be represented as sums of elementary tensors.
This leads to the question of tensor ranks, i.e., the minimum number of simple tensors that can be added to obtain a given tensor.

By construction the map $\otimes$ is bilinear, which written out means
\[
	\begin{array}{rcl}
		(u_1 + u_2) \otimes v & = & u_1 \otimes v + u_2 \otimes v, \\[2pt]
		u \otimes (v_1 + v_2) & = & u \otimes v_1 + u \otimes v_2, \\[2pt]
		(\alpha\, u) \otimes v \;=\; \alpha\,(u \otimes v) & = & u \otimes (\alpha\, v),
	\end{array}
\]
for all $u, u_1, u_2 \in U$, all $v, v_1, v_2 \in V$ and every scalar $\alpha \in K$.
Two consequences are worth recording.
\begin{enumerate}
	\item Because $(\alpha\, u) \otimes v = u \otimes (\alpha\, v)$, the parentheses are unnecessary and we simply write $\alpha\, u \otimes v$.
	\item Expanding a product of two linear combinations gives
	\[
		\Bigl( \sum_{i=1}^{n} \alpha_i\, u_i \Bigr) \otimes \Bigl( \sum_{j=1}^{m} \beta_j\, v_j \Bigr)
		= \sum_{i=1}^{n} \sum_{j=1}^{m} \alpha_i\, \beta_j \; u_i \otimes v_j .
	\]
\end{enumerate}

A final word of caution: the tensor product is in general \textbf{not} commutative.\footnote{Even when $U = V$, one usually has $u \otimes v \neq v \otimes u$, so the order of the two factors must be kept.}

\subsection{Multilinearity, Dual Spaces and Determinant}\label{subsec:rstensors}
Though the tensor product was introduced so that \emph{multilinearity} could be studied with the familiar tools of \emph{linearity}, once the construct is in place, the tensors themselves become the object of interest.
This is especially the case for forming tensor products of copies of primal and dual version of a \emph{single} finite-dimensional space $V$.
For $r, s \ge 0$ this gives the \textbf{tensor space of type $(r,s)$},
\[
	T^{r}_{s}(V) = \underbrace{V \otimes \cdots \otimes V}_{r\ \mathrm{times}} \otimes \underbrace{V^{*} \otimes \cdots \otimes V^{*}}_{s\ \mathrm{times}},
\]
whose elements are the \textbf{$(r,s)$-tensors}, classically called $r$ times \emph{contravariant} and $s$ times \emph{covariant}.

Tensors of type $(r,s)$ can be investigated from different points of view -- quite importantly from the view as mappings selecting some of the axis components as maps from the corresponding (bi-)dual spaces and investigating those mapping properties.
\end{document}
